% Generated by scripts/check_shared.py --emit; illustrative, no medium binding.
\begin{table}[htbp]\centering\small
\begin{tabular}{@{}llrrrrr@{}}\toprule
示意 & 情景 & $\Delta_S$ (ns) & $\Delta_R$ (ns) & $\rho$ (GB/s) & $\tau$ (GB/s) & $\RI^*$\\\midrule
1 & 短时隙 & 1668 & 24 & 0.07674 & 0.6667 & 0.1151\\
1 & 参考 & 3850 & 60 & 0.03325 & 0.2667 & 0.1247\\
1 & 长时隙 & 7700 & 120 & 0.01662 & 0.1333 & 0.1247\\
\midrule
2 & 短时隙 & 1284 & 11286 & 0.09969 & 0.02268 & 4.395\\
2 & 参考 & 3850 & 22590 & 0.03325 & 0.01133 & 2.934\\
2 & 长时隙 & 10260 & 45180 & 0.01248 & 0.005666 & 2.202\\
\bottomrule\end{tabular}
\caption{两个算例的成对情景。时间列保留算术结果用于复算，非器件测量精度；GB 为 $10^9$ Byte。}
\label{tab:examples}\end{table}
示意 1 的展示范围为 $\rho$=0.016--0.077 GB/s、$\tau$=0.13--0.67 GB/s，成对 ridge 为 0.11--0.13；独立端点外包络为 0.024--0.58。
示意 2 的展示范围为 $\rho$=0.012--0.1 GB/s、$\tau$=0.0056--0.023 GB/s，成对 ridge 为 2.2--4.4；独立端点外包络为 0.55--18。
